\documentclass[11pt]{article}

\usepackage[margin=1in]{geometry}
\usepackage[T1]{fontenc}
\usepackage[utf8]{inputenc}
\usepackage{amsmath,amssymb}
\usepackage{booktabs}
\usepackage{enumitem}
\usepackage{microtype}
\usepackage{hyperref}

\emergencystretch=2em

\hypersetup{
  colorlinks=true,
  linkcolor=black,
  citecolor=black,
  urlcolor=blue
}

\title{Free-Energy Selection for Predicate Encapsulation in Sparse Deterministic Solvers}
\author{GKM research manuscript}
\date{\today}

\begin{document}
\maketitle

\begin{abstract}
Open-ended artificial evolution cannot be obtained by optimizing a fixed finite
benchmark alone. With positive complexity pressure, a fixed benchmark eventually
selects compression of the best observable behavior, not continuing innovation.
This manuscript develops a narrower experimental step toward an open-ended
research program: sparse deterministic solvers are selected by a free-energy
objective that trades empirical loss against explicit description length, and
reusable predicate macros are admitted only when their library cost is repaid by
reducing repeated task code. The current result is deliberately limited. The
primitive observations are hand-defined; the system does not yet discover
perceptual atoms from pixels, turtle programs, or raw geometry. What is shown is
that, given such atoms, free-energy accounting selects a shared predicate macro
only when a latent conjunction is reused across tasks or across factored
disjunctive branches, and removes the effect when sharing is made unavailable.
We formalize the objective, report the internal abstraction-emergence control
experiment, relate it to minimum description length, predicate invention,
library learning, automata learning, and Bongard-style concept induction, and
lay out the next step: evolving finite-state predicate recognizers over
symbolic Bongard-LOGO action programs under the same accounting protocol.
\end{abstract}

\section{Introduction}

The working thesis of this repository is that open-ended artificial evolution is
possible only when a local selection rule is coupled to an expanding ecology.
The local rule used here is a free-energy objective,
\begin{equation}
  F_{\lambda}(h;D)=R_D(h)+\lambda C(h),
  \label{eq:basic-free-energy}
\end{equation}
where $h$ is an agent, automaton, transducer, or task solver, $R_D(h)$ is
empirical risk on the current data or environment distribution, $C(h)$ is a
description-length proxy, and $\lambda>0$ controls parsimony pressure. This form
is close in spirit to minimum description length \cite{rissanen1978modeling,
grunwald2007minimum,grunwaldroos2019mdl} and to the loss-complexity
structure-function view in Kolpakov's recent free-energy formulation
\cite{kolpakov2025loss}. It is also compatible with broader free-energy
language in adaptive systems \cite{friston2010free}, although the present work
uses the term operationally: selection minimizes a loss plus a complexity cost.

Equation~\ref{eq:basic-free-energy} is not by itself an open-endedness
principle. On a fixed finite benchmark, every candidate program induces a finite
behavior vector. Once the best attainable vector is reached, extra machinery is
penalized unless it reduces risk on an already specified item. In that regime,
free-energy minimization compresses. Open-endedness requires that the risk term
itself change as agents solve niches, create artifacts, or expose new
frontiers:
\begin{equation}
  Q_{t+1}=G(Q_t,\mathcal{P}_t,\mathcal{A}_t),
  \label{eq:ecology}
\end{equation}
where $Q_t$ is the environment distribution, $\mathcal{P}_t$ is the population,
and $\mathcal{A}_t$ is an archive of earlier agents, environments, and
counterexamples. Under this view, free energy selects locally; the ecology
expands globally.

The immediate research question is smaller and more measurable. Before claiming
open-endedness, we need to show that the complexity term can do real structural
work. In particular, can it favor reusable abstractions only when reuse earns
its cost? The current experiment studies that question with sparse deterministic
rules over predefined symbolic observations. The result is not predicate
invention from raw input. It is predicate \emph{encapsulation}: a repeated
conjunction can be defined once as a macro and called cheaply by several task
rules.

\paragraph{Contribution.}
The present manuscript makes four contributions.
\begin{enumerate}[leftmargin=1.5em]
  \item It states a free-energy objective for sparse deterministic solvers and
  library-bearing task families.
  \item It defines the exact loss and complexity accounting used in the current
  abstraction-emergence experiment.
  \item It reports a controlled internal result showing when a shared predicate
  macro is selected, when it is not selected, and why the no-share ablation
  matters.
  \item It positions the result relative to established work on predicate
  invention, program-library learning, automata learning, and Bongard-style
  visual reasoning, while keeping the novelty claim narrow.
\end{enumerate}

\section{Free-energy selection and structure functions}

\subsection{Risk, complexity, and lambda sweeps}

For a labeled concept task with data $D=\{(x_i,y_i)\}_{i=1}^{n}$ and a
deterministic solver $h:X\to\{0,1\}$, the basic empirical risk is
\begin{equation}
  R_D(h)=\frac{1}{n}\sum_{i=1}^{n}\mathbf{1}[h(x_i)\ne y_i].
  \label{eq:risk}
\end{equation}
For a set of tasks $\mathcal{T}=\{1,\ldots,T\}$, each with its own dataset
$D_t$ and solver $h_t$, the total task risk is
\begin{equation}
  R_{\mathcal{D}}(H)=\sum_{t=1}^{T}R_{D_t}(h_t),
  \qquad H=\{h_t\}_{t=1}^{T}.
  \label{eq:total-risk}
\end{equation}
The experiments sweep $\lambda$ rather than treating it as a single tuned
hyperparameter:
\begin{equation}
  F_t(\lambda)=\inf_{h}\left[R_t(h)+\lambda C(h)\right].
  \label{eq:lambda-frontier}
\end{equation}
This estimates a loss-complexity frontier: the minimum complexity needed to
reach a given risk level. Following the structure-function viewpoint
\cite{kolpakov2025loss}, the frontier itself is an object of study. Sharp
changes in selected complexity across $\lambda$ are diagnostics of structural
tradeoffs, not merely hyperparameter artifacts.

\subsection{The complexity ratchet}

A more complex mutation $h'$ is selected over $h$ only when
\begin{equation}
  R_D(h)-R_D(h')>\lambda\left[C(h')-C(h)\right].
  \label{eq:ratchet}
\end{equation}
This is the complexity ratchet. It does not reward complexity as such; it admits
complexity only when risk reduction pays for it. In an expanding ecology,
previously useless structure can become useful because later tasks make new risk
terms visible. In a fixed ecology, the same rule eventually favors compression.

\section{Sparse deterministic substrates}

The repository currently uses three increasingly abstract substrates.

\paragraph{Grid foraging automata.}
The embodied baseline evolves sparse finite-state automata on grid foraging
tasks. A rule has the form
\begin{equation}
  (\mathrm{state},\mathrm{previous\ move},\mathrm{food\ azimuth})
  \mapsto
  (\mathrm{move\ sequence},\mathrm{next\ state}).
\end{equation}
Undefined inputs halt the episode; there is no random fallback policy. The
default complexity mode charges the full encoded sparse rule set, including
unused carried rules and long move sequences.

\paragraph{Register transducers.}
The pattern substrate evolves deterministic transducers over opaque object
tokens. Token identity is unavailable to the rule key. Primitive tiers expose
streaming, registers, equality comparisons, and bidirectional input motion. This
separates representability questions from search questions: if a one-way stream
cannot express a reversal without unbounded memory, the failure is a substrate
failure, not merely an optimizer failure.

\paragraph{Predicate-library scaffold.}
The current abstraction experiment is a non-visual Bongard-style scaffold. Each
object is represented by a set of primitive boolean observations, such as
\begin{equation}
  \texttt{low\_closure\_error},\quad
  \texttt{high\_hull\_fill},\quad
  \texttt{turn\_balanced},\quad
  \texttt{has\_curve},\quad
  \texttt{thin}.
\end{equation}
These atoms are hand-defined. The experiment asks whether a repeated conjunction
over them is encapsulated as a reusable predicate macro.

\section{Library-aware free energy}

Let $L$ be a library of reusable predicates and $H=\{h_t\}$ be task solvers that
may call predicates in $L$. The library-aware objective is
\begin{equation}
  F_{\lambda}(L,H;\mathcal{D})=
  \sum_{t=1}^{T} R_{D_t}(h_t;L)
  +\lambda\left(C_{\mathrm{lib}}(L)+
  \sum_{t=1}^{T}C_{\mathrm{task}}(h_t\mid L)\right).
  \label{eq:library-free-energy}
\end{equation}
This objective is the central accounting device. A macro is useful only if its
definition cost is smaller than the duplicated inline code it replaces.

\subsection{Rule language}

Task rules are disjunctive normal forms. A solver returns true if at least one
clause is satisfied:
\begin{equation}
  h(x)=\bigvee_{j=1}^{m}\bigwedge_{a\in A_j} a(x).
\end{equation}
The reusable latent conjunction in the current experiment is
\begin{equation}
  A =
  \texttt{low\_closure\_error}
  \wedge
  \texttt{high\_hull\_fill}
  \wedge
  \texttt{turn\_balanced}.
  \label{eq:solid-loop}
\end{equation}
The selected macro is named \texttt{solid\_loop}. Example support tasks include
$A\wedge\texttt{has\_curve}$,
$A\wedge\texttt{thin}$, and
$A\wedge\texttt{symmetric\_hint}$. The key disjunctive diagnostic is
\begin{equation}
  (A\wedge B)\vee(A\wedge C),
  \label{eq:or-factor}
\end{equation}
where $B=\texttt{has\_curve}$ and $C=\texttt{thin}$. The unrelated OR control is
\begin{equation}
  B\vee C.
  \label{eq:or-control}
\end{equation}

\subsection{Complexity accounting}

The current scaffold uses the following costs:
\begin{align}
  C(\mathrm{primitive\ atom\ reference}) &= 1,\\
  C(\mathrm{rule\ overhead}) &= 1,\\
  C(\mathrm{macro\ definition}) &= 1+\left|\mathrm{body}\right|,\\
  C(\mathrm{shared\ macro\ call}) &= \rho,\qquad \rho=0.35.
\end{align}
For a DNF rule $h$, the task complexity is
\begin{equation}
  C_{\mathrm{task}}(h\mid L)
  =
  1+
  \sum_{j=1}^{m}\sum_{\ell\in A_j}
  c(\ell\mid L),
  \label{eq:task-complexity}
\end{equation}
where $c(\ell\mid L)=1$ for primitive atom references and $c(\ell\mid L)=\rho$
for a shared macro call. In the no-share ablation, macro syntax is allowed but a
macro call is charged as
\begin{equation}
  c_{\mathrm{no\ share}}(\ell\mid L)=\rho+C_{\mathrm{def}}(\ell),
\end{equation}
so every use repays the full definition. This ablation distinguishes genuine
shared description-length savings from mere syntactic convenience.

These numbers explain the observed controls. The unrelated OR control has
complexity $1+1+1=3$. The inline factored OR in
Equation~\ref{eq:or-factor} has complexity $1+4+4=9$. The shared version pays
$4$ for defining \texttt{solid\_loop} and
$1+(0.35+1)+(0.35+1)=3.7$ for the task rule, for total complexity $7.7$.

\section{Abstraction-emergence experiment}

\subsection{Protocol}

Each task is generated from a deterministic boolean concept over the primitive
atoms. Positive, validation, and hidden examples are sampled from the finite
object universe, with hard negatives prioritized by overlap with the positive
clauses. Conditions are:
\begin{description}[leftmargin=2.2em,style=nextline]
  \item[inline]
  Each task solves directly over primitive atoms.
  \item[shared]
  A macro may be defined once and called cheaply by several task rules.
  \item[no\_share]
  Macro syntax is available, but the macro definition is charged at every use.
  \item[oracle]
  A privileged task predicate is supplied directly. This is an upper bound, not
  a discovery condition.
\end{description}
The default sweep uses $\lambda\in\{0.005,0.01,0.02,0.04\}$. Selection minimizes
training free energy for each $\lambda$ and then chooses the simplest candidate
among those with best validation loss. Hidden examples are reported after
selection as an audit of generalization within the finite symbolic scaffold.

\subsection{Results}

Table~\ref{tab:abstraction-results} reports the key rows from the May 23, 2026
run. All listed rows reach zero train, validation, and hidden loss; the
selection differences are therefore entirely due to complexity.

\begin{table}[t]
\centering
\small
\begin{tabular}{@{}llrrlp{0.33\linewidth}@{}}
\toprule
Scenario & Condition & Loss & Complexity & Macro & Interpretation \\
\midrule
single & inline & 0.00 & 5.00 & none & One task has no reuse pressure. \\
single & shared & 0.00 & 5.00 & none & Sharing is available but not selected. \\
multi & inline & 0.00 & 15.00 & none & Three support tasks duplicate the same core conjunction. \\
multi & shared & 0.00 & 11.05 & \texttt{solid\_loop} & Defining the repeated core once lowers total description length. \\
multi & no\_share & 0.00 & 15.00 & none & The effect disappears when the definition is repaid at every call. \\
or\_control & shared & 0.00 & 3.00 & none & The simple $B\vee C$ control contains no repeated substructure. \\
or\_factor & inline & 0.00 & 9.00 & none & $(A\wedge B)\vee(A\wedge C)$ repeats $A$ inline. \\
or\_factor & shared & 0.00 & 7.70 & \texttt{solid\_loop} & Factoring the repeated $A$ lowers free energy. \\
or\_factor & no\_share & 0.00 & 9.00 & none & Syntax alone is insufficient without shared cost accounting. \\
or\_factor\_transfer & shared & 0.00 & 2.35 & \texttt{solid\_loop} & A new task using $A$ has lower marginal cost than inline complexity $5.00$. \\
\bottomrule
\end{tabular}
\caption{Key abstraction-emergence controls. The macro body is
\texttt{low\_closure\_error AND high\_hull\_fill AND turn\_balanced}.}
\label{tab:abstraction-results}
\end{table}

\subsection{Interpretation}

The experiment establishes a narrow but useful internal control. Single-task
pressure does not create a macro. An unrelated OR does not create a macro. A
multi-task family with repeated latent structure does create a macro. A
factored disjunction of the form $(A\wedge B)\vee(A\wedge C)$ also creates a
macro. The no-share ablation removes the effect, showing that the selected
macro is a shared-description-length phenomenon.

This is not yet full predicate invention. The atom
\texttt{low\_closure\_error} is not learned. The conjunction in
Equation~\ref{eq:solid-loop} is selected as a reusable macro over a supplied
primitive vocabulary. The next experiment must evolve the predicate recognizer
itself, for example as a finite-state acceptor over LOGO action-program traces,
then ask whether downstream solvers call that acceptor when it reduces
Equation~\ref{eq:library-free-energy}.

\section{Developmental overcapacity}

The broader automata experiments suggest another important mechanism. A useful
behavior may first be discovered in an overbuilt part of the search space and
only later become compressible. This is not a failure of the complexity term.
Equation~\ref{eq:ratchet} still penalizes unused or redundant structure. The
point is that the globally minimal basin may be hard to enter from cold
stochastic mutation. Larger temporary capacity can expose a working behavior;
after that, complexity pressure can favor smaller descendants.

This developmental-overcapacity hypothesis is distinct from the macro result.
In the macro scaffold, all relevant candidates are enumerated or derived from
inline DNF branches. In the sparse Bongard-style classifier experiments, by
contrast, some boundary-equality concepts are found only when the search is
allowed more rules than the final exact solver uses. The manuscript therefore
treats overcapacity as a search-accessibility hypothesis requiring separate
rate experiments, not as proof of abstraction emergence.

\section{Path to Bongard-style predicate discovery}

Bongard problems are a natural target because they require few-shot concept
induction from positive and negative panels \cite{bongard1970pattern}. Modern
Bongard-LOGO provides program-guided synthetic concepts with ground-truth LOGO
generation programs \cite{nie2020bongardlogo}. Recent symbolic-grounding
diagnostics show that symbolic action-program inputs can greatly reduce the
representational bottleneck relative to raw pixels \cite{vaishnav2026symbolic}.

The planned progression is:
\begin{enumerate}[leftmargin=1.5em]
  \item Keep raw perception out of the loop initially. Use symbolic LOGO action
  programs so failures can be attributed to rule induction rather than vision.
  \item Evolve predicate automata that map an action program or structured trace
  to true/false.
  \item Charge each predicate automaton as library complexity.
  \item Let task solvers call learned predicate automata when doing so lowers
  total free energy.
  \item Run the same controls as Table~\ref{tab:abstraction-results}: single
  task, unrelated OR, factored OR, no-share, oracle metadata, and transfer.
\end{enumerate}
This would move from predicate encapsulation over predefined atoms toward
predicate discovery in a controlled symbolic visual-reasoning benchmark.

\section{Related work}

\paragraph{Predicate invention and ILP.}
Predicate invention is a central, long-standing problem in inductive logic
programming, not a new idea introduced here. Muggleton and Buntine
\cite{muggleton1988machine} introduced early predicate invention by inverting
resolution; later work formalized predicate invention and utilization
\cite{muggleton1994predicate} and ILP more broadly
\cite{muggletonderaedt1994ilp}. Modern surveys and systems continue to treat
predicate invention as powerful but difficult
\cite{cropper2021ilp30,cropper2021predicate}. Meta-interpretive learning is
especially relevant because it supports predicate invention in grammatical and
program-induction settings \cite{muggleton2014meta}.

\paragraph{Library learning and compression.}
DreamCoder learns symbolic abstractions and neural search guidance across tasks
\cite{ellis2020dreamcoder}. Stitch reframes abstraction discovery as efficient
top-down synthesis for library learning \cite{bowers2023stitch}, and LILO adds
LLM-assisted synthesis and documentation around compression-based library
construction \cite{grand2023lilo}. The present experiment is close in spirit but
uses a smaller deterministic substrate and explicit library-plus-task
free-energy accounting.

\paragraph{Automata learning.}
Gold's work on language identification \cite{gold1967language} and Angluin's
query-based learning of regular sets \cite{angluin1987learning} are
foundational for thinking about sparse finite-state recognizers. Automata
extraction from recurrent networks \cite{weiss2018extracting} is adjacent but
has a different goal: interpreting an already trained neural model rather than
evolving a compact symbolic solver under description-length pressure.

\paragraph{Temporal abstraction.}
The options framework \cite{sutton1999between}, Option-Critic
\cite{bacon2017option}, and skill chaining \cite{konidaris2009skill} provide a
reinforcement-learning analogy: reusable temporally extended structures can
make later behavior cheaper to express. The current macros are sensory
predicates rather than action policies, but the same accounting question
appears: when does a reusable internal object earn its cost?

\paragraph{Bongard reasoning.}
Classical Bongard problems and later Bongard-LOGO benchmarks
\cite{bongard1970pattern,nie2020bongardlogo} motivate the concept-learning
setting. Depeweg et al. used visual languages and pragmatic reasoning
\cite{depeweg2018solving}; Sonwane et al. combined program synthesis and ILP
for Bongard-style reasoning \cite{sonwane2021using}. Bongard-HOI
\cite{jiang2022bongardhoi} and Bongard-OpenWorld \cite{wu2023bongardopenworld}
are harder downstream targets, but they introduce natural-image and
open-vocabulary burdens before the free-energy predicate story is isolated.

\section{Limitations}

The current evidence is intentionally modest.
\begin{enumerate}[leftmargin=1.5em]
  \item Primitive observations are predefined. The experiment does not discover
  \texttt{low\_closure\_error}, \texttt{high\_hull\_fill}, or
  \texttt{turn\_balanced}.
  \item Candidate macros are enumerated for conjunctive support tasks and
  derived from repeated inline DNF branches for the OR-factor task. This is a
  scaffold for studying accounting, not a complete invention mechanism.
  \item The object universe is finite and symbolic. It is useful for controls
  but too clean to support claims about natural visual abstraction.
  \item The result demonstrates compression under fixed tasks. Open-endedness
  still requires an endogenous environment generator as in Equation~\ref{eq:ecology}.
  \item Macro selection is sensitive to the cost model. That sensitivity is a
  feature for structure-function analysis, but it must be reported by lambda and
  cost sweeps rather than hidden behind a single setting.
\end{enumerate}

\section{Conclusion}

The present result is a small but necessary step. Under explicit free-energy
accounting, a sparse deterministic solver does not create a reusable predicate
macro merely because macro syntax exists. It creates one when repeated structure
makes the shared definition cheaper than duplicated inline rules, and it stops
doing so when sharing is made unavailable. This supports the local selection
part of the larger thesis: complexity can grow or reorganize when it purchases
loss reduction or transfer, and otherwise should be compressed away. The next
scientific step is to replace predefined primitive atoms with evolved
finite-state predicate recognizers on symbolic Bongard-LOGO traces and test
whether the same library/free-energy controls still separate real abstraction
from notation.

\section*{Reproducibility note}

The abstraction-emergence table is reproduced from the repository root with:
\begin{verbatim}
python3 experiments/run_abstraction_emergence.py
python3 experiments/run_abstraction_emergence.py --scenario or_factor --show-rules
python3 -m unittest tests.test_abstraction_emergence
\end{verbatim}

\bibliographystyle{plain}
\bibliography{references}

\end{document}
